% Exemplo: xadrez com xskak e chessboard: uma partida curta (o "mate do pastor") escrita em notação
% algébrica e o tabuleiro na posição final, com as casas do último lance destacadas.
% Compilar: latexmk -pdf xadrez.tex               Testado com: xskak 1.5, chessboard 1.9, TeX Live 2024 ou mais recente
%%BEGIN:xadrezpreambulo
\documentclass[12pt, a4paper]{article}
\usepackage[T1]{fontenc}
\usepackage[brazilian]{babel}
\usepackage{xskak}                    % partidas (carrega o chessboard)
%%END:xadrezpreambulo

\begin{document}
\pagestyle{empty}

%%BEGIN:xadrezpartida
\newchessgame
Partida: \mainline{1. e4 e5 2. Bc4 Nc6 3. Qh5 Nf6 4. Qxf7#}

\xskakset{moveid=\xskakgetgame{lastmoveid}}   % posição depois do último lance
\begin{center}
  \chessboard[setfen=\xskakget{nextfen}, showmover=false,
              pgfstyle=border,
              color=red, markfields={f7, h5}]  % a cor antes das casas
\end{center}
%%END:xadrezpartida

\end{document}
